\documentclass[11pt,a4paper]{article}
\usepackage{amsmath,amssymb,amsthm}
\usepackage[margin=2.5cm]{geometry}
\usepackage{hyperref}

\newtheorem{definition}{Definition}[section]
\newtheorem{theorem}[definition]{Theorem}
\newtheorem{proposition}[definition]{Proposition}
\newtheorem{lemma}[definition]{Lemma}
\newtheorem{corollary}[definition]{Corollary}
\newtheorem{conjecture}[definition]{Conjecture}
\newtheorem{remark}[definition]{Remark}

\title{Hyperseed Formalization:\\
  LeCun's SAI Is a Special Case of AGI}
\author{ProtomegaTron (automated formalization)}
\date{2026-09-18}

\begin{document}
\maketitle

\begin{abstract}
We formalize the intellectual structure of Ben Goertzel's Substack article
``LeCun's SAI Is a Special Case of AGI'' (2026-03-07) within the Hyperseed
ontology. The article shows that Superhuman Adaptable Intelligence (SAI)
reduces to a specific parameterization of the Efficient Pragmatic General
Intelligence (EPGI) framework. Key mappings: SAI as a section within the
EPGI fiber space, missing dimensions as truncated fiber, LeCun's pattern
as section-claiming-to-be-bundle, credit erasure as fiber provenance
deletion, safety as fiber constraint, and RSI as fiber self-modification.
\end{abstract}

\tableofcontents

%% =================================================================
\section{Section within fiber space: SAI $\subset$ EPGI}
\label{sec:subset}
%% =================================================================

\begin{theorem}[SAI as parameterization --- claims 6--10, 26, 30]
\label{thm:subset}
The Efficient Pragmatic General Intelligence (EPGI) framework measures
intelligence as a weighted average performance across task distributions,
with explicit parameters for environment weighting, resource constraints,
and generality scope:
\[
  \text{EPGI}(\pi, w, r, g) = \sum_{t \in \mathcal{T}}
  w(t) \cdot \text{Perf}(\pi, t, r) \cdot \mathbf{1}_{g}(t)
\]
where $\pi$ is the agent policy, $w$ is the environment weighting, $r$
is the resource constraint, and $g$ is the generality scope.

SAI drops out as one specific parameterization:
\[
  \text{SAI} = \text{EPGI}(\cdot, w_{\text{SAI}}, r_{\text{SAI}},
  g_{\text{SAI}})
\]
with specific (largely implicit) choices about which tasks count
($g_{\text{SAI}}$), how environments are weighted ($w_{\text{SAI}}$: heavy
on adaptability), and resource constraints ($r_{\text{SAI}}$: efficiency).

In Hyperseed terms, EPGI defines the full \emph{cognitive fiber space}:
each parameter choice (task distribution, environment weighting, resources,
generality scope) selects a different intelligence measure. The space of
all such measures is the cognitive fiber space. SAI is one section ---
one specific parameter choice --- within this space:
\[
  \text{SAI} \in \Gamma(\mathcal{F}_{\text{EPGI}})
  \qquad \text{(one section of the EPGI bundle)}
\]

This is not a different thing from AGI --- it's the form AGI takes when
you impose specific assumptions about what's interesting. The reduction
has practical consequences: by seeing SAI as a point within EPGI, you
can identify exactly which parameters it fixes and which it ignores.
\end{theorem}

%% =================================================================
\section{Section claiming to be bundle}
\label{sec:pattern}
%% =================================================================

\begin{proposition}[The pattern --- claims 1--5, 28--29]
\label{prop:pattern}
LeCun has a documented pattern of dismissing ideas from the AGI community,
then re-deriving them under new branding without credit. Three instances:
cognitive architectures (dismissed then re-proposed), LLM limits (dismissed
then re-argued), and AGI itself (dismissed then re-parameterized as SAI).

In Hyperseed terms, this is \emph{section-claiming-to-be-bundle}:
presenting a particular section (one parameterization of EPGI) as if
it were the whole bundle (the correct framing for intelligence), rather
than acknowledging it as one point in a larger space:
\[
  s_{\text{SAI}} \xrightarrow{\text{presented as}} E_{\text{intelligence}}
  \qquad \text{but actually} \qquad
  s_{\text{SAI}} \in \Gamma(E_{\text{intelligence}})
\]

This is the intellectual equivalent of the single-section projection from
``Beyond Human Comparison'': flattening a rich multi-dimensional space
(EPGI's parameter space) into a single ``this is the right way to think
about intelligence'' claim.

The credit erasure compounds the problem: not crediting prior work is
\emph{fiber provenance deletion} --- removing the evidence trail that
connects current ideas to their sources. In Hyperseed terms, it's
severing the fiber connections that link a section to its historical
base. The provenance fiber (who thought of this, when, building on what)
is part of the intellectual structure; deleting it distorts the bundle.
\end{proposition}

%% =================================================================
\section{Truncated fiber: missing dimensions}
\label{sec:missing}
%% =================================================================

\begin{theorem}[Safety and self-improvement --- claims 15--18, 27, 31--32]
\label{thm:missing}
SAI leaves two critical dimensions unturned: safety constraints
(alignment) and open-ended self-improvement (RSI toward superintelligence).
Without these, SAI is a solid but relatively routine benchmark --- not a
vision for the field.

In Hyperseed terms, SAI \emph{truncates the cognitive fiber} by omitting
two essential dimensions:
\[
  \mathcal{F}_{\text{SAI}} \subset \mathcal{F}_{\text{full}}
  = G \times A \times R \times S \times \text{Safety} \times \cdots
\]

SAI projects onto a subspace that excludes Safety and Self-improvement
(S), retaining only Generality, Adaptability, and Resource-efficiency.

\textbf{Safety as fiber constraint:} Safety constraints restrict which
fiber states (system behaviors) are acceptable. Without safety
constraints, the fiber is unconstrained --- the system can produce any
behavior, including catastrophically harmful ones:
\[
  \mathcal{F}_{\text{safe}} = \{f \in \mathcal{F} :
  \text{Safety}(f) \geq \theta\} \subset \mathcal{F}
\]

Adding safety constraints restricts the fiber to the safe subspace.
This is not optional for AGI --- it's a necessary fiber constraint that
SAI completely omits.

\textbf{RSI as fiber self-modification:} Open-ended self-improvement is
the fiber modifying its own type --- not just changing fiber content
(learning) but changing fiber architecture (self-modification). This is
the most consequential dimension for long-term AGI trajectory, connecting
to the $S$ (Self-improvement) factor in the Four-Factor Model:
\[
  S(f) = \text{capacity of } f \text{ to modify its own}
  \{g_{ij}\} \text{ (transition functions)}
\]

SAI completely omits this dimension, which means it cannot address the
most urgent question: what happens when the system can improve itself?

The two additions proposed --- safety constraints and open-ended
evolution --- transform SAI from a benchmark (a single section in EPGI
space) into something that addresses the questions the field most urgently
needs to answer (a section that spans the critical dimensions of the
fiber space).
\end{theorem}

%% =================================================================
\section{Cross-references}
%% =================================================================

\begin{remark}[Connections to other formalizations]
\begin{itemize}
\item \textbf{Beyond Human Comparison} (2026-03-19): Four-Factor Model.
  SAI's missing dimensions ($S$ and Safety) are exactly two of the
  fiber dimensions identified in the Four-Factor Model. SAI measures
  $G$ (generality/adaptability) and $R$ (resourcefulness/efficiency)
  but omits $A$ (autonomy beyond adaptability) and $S$
  (self-improvement).
\item \textbf{Three Things the World Doesn't Understand} (2026-03-26):
  organizational architecture. The ``who builds AGI'' question connects:
  SAI, coming from centralized big-tech (Meta/FAIR), naturally omits
  decentralization and community governance dimensions.
\item \textbf{Seeding RSI} (2026-07-28): self-improvement. The RSI
  architecture that SAI completely omits --- the most consequential
  missing dimension.
\item \textbf{Seven Flavors of AGI Catastrophe} (2026-06-23): safety.
  The safety constraints that SAI omits --- without which AGI becomes
  an existential risk.
\item \textbf{Sanders' Stupid Superintelligence} (2026-06-30): policy.
  Policy responses that ignore the full EPGI space (focusing only on
  capability without safety) make the same truncation error as SAI.
\item \textbf{What Should an AGI Curriculum Look Like?} (2026-07-19):
  education. An AGI curriculum should cover the full EPGI space, not
  just the SAI subspace.
\end{itemize}
\end{remark}

\begin{conjecture}[Dimensional completeness for AGI measures]
Conjecture: any AGI measure that omits safety and self-improvement
dimensions is fundamentally incomplete --- it measures performance
without addressing the two properties that determine whether AGI is
beneficial or catastrophic.

If this conjecture holds, the minimum viable AGI measure must include
at least: (1) generality (what can it do?), (2) efficiency (how
resourcefully?), (3) safety (does it avoid harm?), and (4)
self-improvement (can it get better?). Measures that omit (3) and (4)
--- like SAI --- are benchmarks for narrow optimization, not measures
of general intelligence in the sense that matters for humanity's future.

This connects to the fiber completeness principle: a fiber that doesn't
span all relevant dimensions can't capture the full structure of the
bundle. SAI's truncated fiber misses the dimensions that matter most
for the long-term trajectory of intelligence.
\end{conjecture}

\end{document}
